%%%%%%%%%%%%%%%%%%%%%%%%%%%%%%%%%%%%%%%%%%%%%%%%%%%%%%%%%%%%%%%%%%%%%%%%%%%%%%%
% Chapter 2: Mathematical Background
%%%%%%%%%%%%%%%%%%%%%%%%%%%%%%%%%%%%%%%%%%%%%%%%%%%%%%%%%%%%%%%%%%%%%%%%%%%%%%%

\chapter{Mathematical Background}
\label{ch:background}

This chapter presents the mathematical foundations of Graph Neural Networks required for understanding the models used in this research. We follow the notation and formulations from established literature \cite{kipf2017, kazemi2020, scarselli2008}.

\section{Graph Representation}

\subsection{Basic Definitions}

An \textbf{attributed graph} is represented as a tuple $G = (V, \mat{A}, \mat{X})$ where:

\begin{itemize}
    \item $V = \{v_1, v_2, \ldots, v_n\}$ is the set of vertices (nodes)
    \item $n = |V|$ denotes the number of nodes
    \item $\mat{A} \in \R^{n \times n}$ is the adjacency matrix
    \item $\mat{X} \in \R^{n \times d}$ is the feature matrix where $\mat{X}_i$ represents the $d$-dimensional features of node $v_i$
\end{itemize}

The adjacency matrix $\mat{A}$ encodes the graph structure:
\begin{equation}
    \mat{A}_{i,j} = \begin{cases}
        w_{ij} & \text{if there exists an edge between } v_i \text{ and } v_j \\
        0 & \text{otherwise}
    \end{cases}
\end{equation}

where $w_{ij} \in \R_+$ represents the edge weight.

\subsection{Graph Properties}

\textbf{Undirected graphs:} The adjacency matrix is symmetric, i.e., $\mat{A} = \mat{A}^\top$.

\textbf{Weighted graphs:} Edge weights can represent various properties such as:
\begin{itemize}
    \item \textbf{Distance:} Euclidean distance between nodes (as in our vehicle graphs)
    \item \textbf{Similarity:} Cosine similarity between node features
    \item \textbf{Connection strength:} Normalized affinity scores
\end{itemize}

In our application, each frame produces an \textbf{undirected weighted graph} where edge weights represent Euclidean distances between vehicle centroids.

\section{Message Passing Framework}

Graph Neural Networks operate through a \textbf{message passing} paradigm where node representations are iteratively updated by aggregating information from neighboring nodes.

\subsection{General Message Passing}

The general message passing framework consists of three operations at each layer $l$:

\begin{enumerate}
    \item \textbf{Message computation:} For each edge $(i, j)$, compute a message:
    \begin{equation}
        \vect{m}_{ij}^{(l)} = \phi^{(l)}\left(\vect{h}_i^{(l-1)}, \vect{h}_j^{(l-1)}, e_{ij}\right)
    \end{equation}
    
    \item \textbf{Aggregation:} Aggregate messages from all neighbors:
    \begin{equation}
        \vect{a}_i^{(l)} = \bigoplus_{j \in \mathcal{N}(i)} \vect{m}_{ij}^{(l)}
    \end{equation}
    where $\mathcal{N}(i)$ denotes the neighborhood of node $i$ and $\bigoplus$ is a permutation-invariant aggregation function (e.g., sum, mean, max).
    
    \item \textbf{Update:} Update node representation:
    \begin{equation}
        \vect{h}_i^{(l)} = \psi^{(l)}\left(\vect{h}_i^{(l-1)}, \vect{a}_i^{(l)}\right)
    \end{equation}
\end{enumerate}

Figure~\ref{fig:message_passing} illustrates the message passing process.

\begin{figure}[H]
    \centering
    % Message Passing Diagram
\begin{tikzpicture}[
    node distance=1.5cm,
    vertex/.style={
        circle,
        draw=black,
        thick,
        minimum size=0.8cm,
        font=\small
    },
    center/.style={vertex, fill=orange!40},
    neighbor/.style={vertex, fill=green!30},
    other/.style={vertex, fill=gray!20},
    arrow/.style={-Stealth, thick},
    messagearrow/.style={-Stealth, thick, blue!70, dashed}
]

% Title
\node[font=\bfseries] at (3, 3.5) {Message Passing in GNN};

% Center node
\node[center] (v) at (3, 1.5) {$v$};

% Neighbor nodes
\node[neighbor] (n1) at (1, 2.5) {$u_1$};
\node[neighbor] (n2) at (2, 3) {$u_2$};
\node[neighbor] (n3) at (4, 3) {$u_3$};
\node[neighbor] (n4) at (5, 2.5) {$u_4$};
\node[neighbor] (n5) at (4.5, 0.5) {$u_5$};
\node[neighbor] (n6) at (1.5, 0.5) {$u_6$};

% Edges
\draw[thick] (v) -- (n1);
\draw[thick] (v) -- (n2);
\draw[thick] (v) -- (n3);
\draw[thick] (v) -- (n4);
\draw[thick] (v) -- (n5);
\draw[thick] (v) -- (n6);

% Message arrows
\draw[messagearrow] (n1) -- (v);
\draw[messagearrow] (n2) -- (v);
\draw[messagearrow] (n3) -- (v);
\draw[messagearrow] (n4) -- (v);
\draw[messagearrow] (n5) -- (v);
\draw[messagearrow] (n6) -- (v);

% Equations
\node[align=left, font=\small] at (8.5, 2.5) {
    \textbf{1. Message:}\\[0.2cm]
    $\vect{m}_{u \to v} = \phi(\vect{h}_u, \vect{h}_v, e_{uv})$\\[0.4cm]
    \textbf{2. Aggregate:}\\[0.2cm]
    $\vect{a}_v = \bigoplus_{u \in \mathcal{N}(v)} \vect{m}_{u \to v}$\\[0.4cm]
    \textbf{3. Update:}\\[0.2cm]
    $\vect{h}_v^{(l)} = \psi(\vect{h}_v^{(l-1)}, \vect{a}_v)$
};

% Legend
\node[center, minimum size=0.5cm] at (0.5, -0.5) {};
\node[font=\tiny, right] at (0.8, -0.5) {Target node};
\node[neighbor, minimum size=0.5cm] at (3, -0.5) {};
\node[font=\tiny, right] at (3.3, -0.5) {Neighbors $\mathcal{N}(v)$};

\end{tikzpicture}





    \caption{Message passing in Graph Neural Networks. Each node aggregates features from its neighbors to update its representation.}
    \label{fig:message_passing}
\end{figure}

\section{Graph Convolutional Networks (GCN)}

The Graph Convolutional Network (GCN) \cite{kipf2017} is a specific instantiation of the message passing framework that has proven highly effective for graph-structured data.

\subsection{GCN Layer Formulation}

A single GCN layer transforms node features as follows:

\begin{equation}
\boxed{
    \mat{Z}^{(l)} = \sigma\left(\tilde{\mat{D}}^{-\frac{1}{2}} \tilde{\mat{A}} \tilde{\mat{D}}^{-\frac{1}{2}} \mat{Z}^{(l-1)} \mat{W}^{(l)}\right)
}
\label{eq:gcn_layer}
\end{equation}

where:
\begin{itemize}
    \item $\tilde{\mat{A}} = \mat{A} + \mat{I}_n$ is the adjacency matrix with added self-loops
    \item $\tilde{\mat{D}}$ is the degree matrix of $\tilde{\mat{A}}$ with $\tilde{\mat{D}}_{ii} = \sum_j \tilde{\mat{A}}_{ij}$
    \item $\mat{W}^{(l)} \in \R^{d_{l-1} \times d_l}$ is the learnable weight matrix for layer $l$
    \item $\sigma(\cdot)$ is a non-linear activation function (typically ReLU)
    \item $\mat{Z}^{(l-1)} \in \R^{n \times d_{l-1}}$ are the input node features
    \item $\mat{Z}^{(l)} \in \R^{n \times d_l}$ are the output node features
\end{itemize}

\subsection{Normalized Adjacency Matrix}

The symmetric normalization $\hat{\mat{A}} = \tilde{\mat{D}}^{-\frac{1}{2}} \tilde{\mat{A}} \tilde{\mat{D}}^{-\frac{1}{2}}$ ensures:

\begin{enumerate}
    \item \textbf{Eigenvalue bounds:} Eigenvalues are bounded in $[-1, 1]$
    \item \textbf{Numerical stability:} Prevents exploding/vanishing gradients
    \item \textbf{Scale invariance:} Feature aggregation is invariant to node degree
\end{enumerate}

The element-wise formula for the normalized adjacency is:
\begin{equation}
    \hat{\mat{A}}_{ij} = \frac{\tilde{\mat{A}}_{ij}}{\sqrt{\tilde{\mat{D}}_{ii}} \cdot \sqrt{\tilde{\mat{D}}_{jj}}}
\end{equation}

\subsection{Per-Node GCN Update}

For a single node $v_i$, the GCN update can be written as:
\begin{equation}
    \vect{h}_i^{(l)} = \sigma\left(\mat{W}^{(l)} \cdot \sum_{j \in \mathcal{N}(i) \cup \{i\}} \frac{1}{\sqrt{d_i + 1} \cdot \sqrt{d_j + 1}} \vect{h}_j^{(l-1)}\right)
\end{equation}

where $d_i$ and $d_j$ are the degrees of nodes $i$ and $j$ respectively.

Figure~\ref{fig:gcn_operation} illustrates the GCN convolution operation.

\begin{figure}[H]
    \centering
    % GCN Operation Diagram
\begin{tikzpicture}[
    node distance=1.2cm,
    box/.style={
        rectangle,
        rounded corners=3pt,
        minimum width=1.6cm,
        minimum height=1.4cm,
        align=center,
        draw=black,
        thick,
        font=\small
    },
    opbox/.style={
        rectangle,
        rounded corners=2pt,
        minimum width=0.8cm,
        minimum height=0.6cm,
        align=center,
        draw=black,
        fill=yellow!30,
        font=\small
    },
    arrow/.style={-Stealth, thick}
]

% Title
\node[font=\bfseries] at (7, 4.5) {Graph Convolution Operation};

% Input features Z^(l-1)
\node[box, fill=cyan!20] (X) at (0, 2) {$\mat{Z}^{(l-1)}$};
\node[font=\scriptsize] at (0, 0.8) {$[n \times d_{in}]$};

% Adjacency A_hat
\node[box, fill=green!20] (A) at (3, 2) {$\hat{\mat{A}}$};
\node[font=\scriptsize] at (3, 0.8) {$[n \times n]$};

% Weight matrix W
\node[box, fill=orange!20] (W) at (6, 2) {$\mat{W}^{(l)}$};
\node[font=\scriptsize] at (6, 0.8) {$[d_{in} \times d_{out}]$};

% Output Z^(l)
\node[box, fill=red!20] (Z) at (12, 2) {$\mat{Z}^{(l)}$};
\node[font=\scriptsize] at (12, 0.8) {$[n \times d_{out}]$};

% Operators
\node[opbox] (mult1) at (1.5, 2) {$\times$};
\node[opbox] (mult2) at (4.5, 2) {$\times$};
\node[opbox, fill=purple!25] (sigma) at (9, 2) {$\sigma(\cdot)$};

% Arrows
\draw[arrow] (X) -- (mult1);
\draw[arrow] (A) -- (mult1);
\draw[arrow] (mult1) -- (A);
\draw[arrow] (A) -- (mult2);
\draw[arrow] (W) -- (mult2);
\draw[arrow] (mult2) -- (sigma);
\draw[arrow] (sigma) -- (Z);

% Intermediate labels
\node[font=\scriptsize, fill=white] at (4.5, 3.2) {$\hat{\mat{A}}\mat{Z}^{(l-1)}\mat{W}^{(l)}$};

% Equation
\node[font=\normalsize, align=center] at (6, -0.8) {
    $\mat{Z}^{(l)} = \sigma\left(\hat{\mat{A}} \mat{Z}^{(l-1)} \mat{W}^{(l)}\right)$
};
\node[font=\small] at (6, -1.6) {
    where $\hat{\mat{A}} = \tilde{\mat{D}}^{-\frac{1}{2}} \tilde{\mat{A}} \tilde{\mat{D}}^{-\frac{1}{2}}$
};

\end{tikzpicture}

    \caption{Graph Convolution operation. Node features are aggregated from neighbors with degree-normalized weights, then transformed through a learnable weight matrix.}
    \label{fig:gcn_operation}
\end{figure}

\section{Dynamic Graph Neural Networks}

Dynamic graphs extend static graphs by incorporating temporal evolution. Following \cite{kazemi2020}, we distinguish two types of dynamic graphs.

\subsection{Discrete-Time Dynamic Graphs (DTDG)}

A DTDG represents the graph as a sequence of snapshots:
\begin{equation}
    \mathcal{G} = \left[G^{(1)}, G^{(2)}, \ldots, G^{(T)}\right]
\end{equation}

where each $G^{(t)} = (V^{(t)}, \mat{A}^{(t)}, \mat{X}^{(t)})$ is a graph at time step $t$.

\textbf{Properties of DTDGs:}
\begin{itemize}
    \item Nodes may appear or disappear across time steps
    \item Edge structure may change between snapshots
    \item Node features may evolve over time
\end{itemize}

Our video anomaly detection problem naturally forms a DTDG where each frame corresponds to a time step.

\subsection{Continuous-Time Dynamic Graphs (CTDG)}

A CTDG represents the graph as a stream of timestamped events:
\begin{equation}
    \mathcal{G} = \{(u_i, v_i, t_i, \vect{e}_i)\}_{i=1}^{N}
\end{equation}

where each event represents an interaction between nodes $u_i$ and $v_i$ at time $t_i$ with optional features $\vect{e}_i$.

Figure~\ref{fig:dynamic_graph_types} illustrates both dynamic graph representations.

\begin{figure}[H]
    \centering
    % Dynamic Graph Types Diagram
\begin{tikzpicture}[
    node distance=0.6cm,
    vertex/.style={
        circle,
        draw=black,
        thick,
        minimum size=0.5cm,
        font=\tiny
    },
    present/.style={vertex, fill=green!40},
    absent/.style={vertex, fill=gray!20, dashed},
    arrow/.style={-Stealth, thick}
]

% Title
\node[font=\bfseries] at (6, 5.5) {Dynamic Graph Representations};

% ===== DTDG Section =====
\node[font=\small\bfseries] at (2, 4.5) {(a) DTDG};

% Time labels
\node[font=\tiny] at (0.5, 3.5) {$t=1$};
\node[font=\tiny] at (2.5, 3.5) {$t=2$};
\node[font=\tiny] at (4.5, 3.5) {$t=3$};

% Snapshot 1
\node[present] (a1) at (0, 2.5) {};
\node[present] (b1) at (1, 2.5) {};
\node[present] (c1) at (0.5, 1.5) {};
\draw[thick] (a1) -- (b1);
\draw[thick] (a1) -- (c1);
\draw[thick] (b1) -- (c1);

% Snapshot 2
\node[present] (a2) at (2, 2.5) {};
\node[present] (b2) at (3, 2.5) {};
\node[present] (c2) at (2.5, 1.5) {};
\node[present] (d2) at (3, 1.5) {};
\draw[thick] (a2) -- (b2);
\draw[thick] (a2) -- (c2);
\draw[thick] (b2) -- (d2);
\draw[thick] (c2) -- (d2);

% Snapshot 3
\node[present] (a3) at (4, 2.5) {};
\node[present] (b3) at (5, 2.5) {};
\node[present] (d3) at (5, 1.5) {};
\draw[thick] (a3) -- (b3);
\draw[thick] (b3) -- (d3);

% Arrows between snapshots
\draw[arrow, gray, dashed] (1.2, 2) -- (1.8, 2);
\draw[arrow, gray, dashed] (3.2, 2) -- (3.8, 2);

% DTDG equation
\node[font=\small] at (2.5, 0.5) {$\mathcal{G} = [G^{(1)}, G^{(2)}, \ldots, G^{(T)}]$};

% ===== CTDG Section =====
\node[font=\small\bfseries] at (9.5, 4.5) {(b) CTDG};

% Timeline
\draw[thick, ->] (7, 2) -- (12, 2);
\node[font=\tiny] at (12, 1.7) {time};

% Events
\fill[blue!60] (7.5, 2) circle (3pt);
\node[font=\tiny, above] at (7.5, 2.2) {$e_1$};
\node[font=\tiny, below] at (7.5, 1.7) {$t_1$};

\fill[red!60] (8.5, 2) circle (3pt);
\node[font=\tiny, above] at (8.5, 2.2) {$e_2$};
\node[font=\tiny, below] at (8.5, 1.7) {$t_2$};

\fill[green!60] (9.5, 2) circle (3pt);
\node[font=\tiny, above] at (9.5, 2.2) {$e_3$};
\node[font=\tiny, below] at (9.5, 1.7) {$t_3$};

\fill[orange!60] (10.5, 2) circle (3pt);
\node[font=\tiny, above] at (10.5, 2.2) {$e_4$};
\node[font=\tiny, below] at (10.5, 1.7) {$t_4$};

\fill[purple!60] (11.5, 2) circle (3pt);
\node[font=\tiny, above] at (11.5, 2.2) {$e_5$};
\node[font=\tiny, below] at (11.5, 1.7) {$t_5$};

% Event descriptions
\node[font=\tiny, align=left] at (9.5, 3.3) {Events: $(u, v, t, e)$};

% CTDG equation
\node[font=\small] at (9.5, 0.5) {$\mathcal{G} = \{(u_i, v_i, t_i, e_i)\}_{i=1}^{N}$};

% Comparison note
\node[font=\tiny, align=center, fill=yellow!20, rounded corners] at (6, -0.5) {
    Our approach uses \textbf{DTDG} with frames as snapshots
};

\end{tikzpicture}

    \caption{Dynamic graph representations: (a) Discrete-Time Dynamic Graph (DTDG) as a sequence of snapshots, (b) Continuous-Time Dynamic Graph (CTDG) as a stream of events.}
    \label{fig:dynamic_graph_types}
\end{figure}

\section{Temporal Modeling with Recurrent Networks}

For processing temporal sequences of graph embeddings, we employ recurrent neural networks, specifically Gated Recurrent Units (GRU) \cite{cho2014}.

\subsection{GRU Formulation}

The GRU updates hidden states using gating mechanisms:

\begin{align}
    \vect{z}_t &= \sigma\left(\mat{W}_z \vect{x}_t + \mat{U}_z \vect{h}_{t-1} + \vect{b}_z\right) & \text{(Update gate)} \\
    \vect{r}_t &= \sigma\left(\mat{W}_r \vect{x}_t + \mat{U}_r \vect{h}_{t-1} + \vect{b}_r\right) & \text{(Reset gate)} \\
    \tilde{\vect{h}}_t &= \tanh\left(\mat{W}_h \vect{x}_t + \mat{U}_h (\vect{r}_t \odot \vect{h}_{t-1}) + \vect{b}_h\right) & \text{(Candidate)} \\
    \vect{h}_t &= (1 - \vect{z}_t) \odot \vect{h}_{t-1} + \vect{z}_t \odot \tilde{\vect{h}}_t & \text{(Output)}
\end{align}

where $\odot$ denotes element-wise multiplication.

\subsection{Combining GCN with GRU}

For dynamic graphs, we combine spatial (GCN) and temporal (GRU) processing:

\begin{enumerate}
    \item \textbf{Spatial encoding:} Apply GCN layers to each snapshot $G^{(t)}$:
    \begin{equation}
        \mat{H}^{(t)} = \text{GCN}\left(G^{(t)}\right) \in \R^{n \times d}
    \end{equation}
    
    \item \textbf{Graph-level pooling:} Aggregate node features to graph-level:
    \begin{equation}
        \vect{g}^{(t)} = \text{POOL}\left(\mat{H}^{(t)}\right) \in \R^{d}
    \end{equation}
    
    \item \textbf{Temporal encoding:} Process sequence with GRU:
    \begin{equation}
        \vect{o} = \text{GRU}\left(\left[\vect{g}^{(1)}, \vect{g}^{(2)}, \ldots, \vect{g}^{(T)}\right]\right)
    \end{equation}
\end{enumerate}

\section{Graph Pooling Operations}

Graph pooling aggregates node features: Mean Pool $\vect{g} = \frac{1}{|V|} \sum_{v \in V} \vect{h}_v$, Max Pool $\vect{g} = \max_{v \in V} \vect{h}_v$, or Sum Pool $\vect{g} = \sum_{v \in V} \vect{h}_v$. Our model uses mean pooling.

\section{Focal Loss for Class Imbalance}

For handling class imbalance in binary classification, we employ Focal Loss \cite{lin2017}.

\subsection{Focal Loss Formulation}

The Focal Loss is defined as:
\begin{equation}
\boxed{
    \text{FL}(p_t) = -\alpha_t (1 - p_t)^\gamma \log(p_t)
}
\label{eq:focal_loss}
\end{equation}

where:
\begin{itemize}
    \item $p_t$ is the predicted probability for the true class
    \item $\gamma \geq 0$ is the focusing parameter
    \item $\alpha_t \in [0, 1]$ is the class balancing factor
\end{itemize}

\subsection{Class Balancing Factor}

The balancing factor $\alpha$ is computed from the training set distribution:
\begin{equation}
    \alpha = \frac{N_{\text{normal}}}{N_{\text{normal}} + N_{\text{anomalous}}}
\end{equation}

For a sample with label $y$:
\begin{equation}
    \alpha_t = \begin{cases}
        \alpha & \text{if } y = 0 \text{ (Normal)} \\
        1 - \alpha & \text{if } y = 1 \text{ (Anomalous)}
    \end{cases}
\end{equation}

\subsection{Focusing Effect}

The modulating factor $(1 - p_t)^\gamma$ has the following properties:
\begin{itemize}
    \item When $p_t \to 1$ (well-classified): $(1 - p_t)^\gamma \to 0$, reducing loss contribution
    \item When $p_t \to 0$ (misclassified): $(1 - p_t)^\gamma \to 1$, maintaining full loss
\end{itemize}

This causes the model to focus on hard-to-classify examples while down-weighting easy examples.

\section{Summary of Mathematical Notation}

Table~\ref{tab:notation} summarizes the mathematical notation used throughout this report.

\begin{table}[H]
\centering
\caption{Summary of Mathematical Notation}
\label{tab:notation}
\begin{tabular}{cl}
\toprule
\textbf{Symbol} & \textbf{Description} \\
\midrule
$G = (V, \mat{A}, \mat{X})$ & Attributed graph \\
$V$ & Set of vertices (nodes) \\
$n = |V|$ & Number of nodes \\
$\mat{A} \in \R^{n \times n}$ & Adjacency matrix \\
$\mat{X} \in \R^{n \times d}$ & Node feature matrix \\
$\mat{W}^{(l)}$ & Weight matrix for layer $l$ \\
$\mat{Z}^{(l)}$ & Node embeddings at layer $l$ \\
$\vect{h}_i$ & Feature vector of node $i$ \\
$\mathcal{N}(i)$ & Neighborhood of node $i$ \\
$\tilde{\mat{A}}$ & Adjacency with self-loops ($\mat{A} + \mat{I}$) \\
$\tilde{\mat{D}}$ & Degree matrix of $\tilde{\mat{A}}$ \\
$\sigma(\cdot)$ & Activation function \\
$\odot$ & Element-wise (Hadamard) product \\
$T$ & Number of time steps \\
$\gamma$ & Focal loss focusing parameter \\
$\alpha$ & Focal loss balancing factor \\
\bottomrule
\end{tabular}
\end{table}




